\section*{Motivacion inicial}

La volatilidad implicita es una de las magnitudes mas usadas para resumir el precio de las opciones. En la practica profesional no se interpreta solamente como un parametro tecnico de un modelo de valoracion, sino como una forma compacta de comparar contratos con strikes, vencimientos y condiciones de mercado distintas. Sin embargo, esa misma utilidad la convierte en una variable delicada: no se observa de manera directa, depende del modelo de inversion empleado y cambia con la microestructura de negociacion.

El proyecto que se documenta en esta memoria parte de operaciones reales de opciones y futuros sobre el IBEX. A partir de contratos, ejecuciones y curvas de tipos, se construye una base de volatilidad implicita, se entrena un conjunto de modelos de regresion y se exponen sus predicciones mediante artefactos explicables. La pregunta central es si se puede construir una herramienta que aproxime la volatilidad implicita con calidad suficiente y, al mismo tiempo, permita inspeccionar por que predice lo que predice.

Esta introduccion se desarrolla con mas detalle en el capitulo \ref{chap:introduccion}. Aqui se anticipa la idea principal: en un contexto financiero regulado, un modelo preciso pero opaco resulta insuficiente. La explicabilidad no se trata como un complemento estetico, sino como un requisito para diagnosticar comportamiento, detectar regiones de riesgo y comunicar limitaciones.
